% ============================================================
% preamble.tex — Enabling the Frontier
% ============================================================

\usepackage[T1]{fontenc}
\usepackage[utf8]{inputenc}
\usepackage{lmodern}
\usepackage{microtype}

% Page geometry
\usepackage[
  top=1.25in,
  bottom=1.25in,
  inner=1.25in,
  outer=1.0in,
  headheight=14pt
]{geometry}

% Typography
\usepackage{setspace}
\setstretch{1.15}
\usepackage{parskip}
\setlength{\parskip}{0.6em}
\setlength{\parindent}{0pt}

% Headers & footers
\usepackage{fancyhdr}
\pagestyle{fancy}
\fancyhf{}
\fancyhead[LE]{\small\textit{Enabling the Frontier}}
\fancyhead[RO]{\small\textit{\leftmark}}
\fancyfoot[C]{\small\thepage}
\renewcommand{\headrulewidth}{0.3pt}
% The book calls its top-level divisions "Sections" (styled that way on
% chapter pages) — keep the running header consistent with that vocabulary.
% Guarded: outline.tex uses the article class, which has no \chaptermark.
\makeatletter
\@ifundefined{chaptermark}{}{%
  \renewcommand{\chaptermark}[1]{\markboth{Section \thechapter.\ #1}{}}%
}
\makeatother

% Colors
\usepackage[dvipsnames]{xcolor}
\definecolor{principleblue}{RGB}{20, 60, 110}
\definecolor{accentgold}{RGB}{180, 140, 40}
\definecolor{lightgray}{RGB}{245, 245, 245}
\definecolor{midgray}{RGB}{120, 120, 120}
\definecolor{tbdred}{RGB}{158, 58, 46}

% Section formatting
\usepackage{titlesec}
\titleformat{\chapter}[display]
  {\normalfont\huge\bfseries\color{principleblue}}
  {\normalfont\large\color{midgray}Section \thechapter}
  {8pt}
  {\Huge}
\titlespacing{\chapter}{0pt}{-10pt}{30pt}

\titleformat{\section}
  {\normalfont\Large\bfseries\color{principleblue}}
  {\thesection}
  {1em}
  {}

\titleformat{\subsection}
  {\normalfont\normalsize\bfseries\color{midgray}\itshape}
  {}
  {0em}
  {}

% Principle box environment
\usepackage{mdframed}
\newmdenv[
  backgroundcolor=lightgray,
  linecolor=principleblue,
  linewidth=2pt,
  topline=false,
  bottomline=false,
  rightline=false,
  innerleftmargin=14pt,
  innerrightmargin=14pt,
  innertopmargin=10pt,
  innerbottommargin=10pt,
  skipabove=12pt,
  skipbelow=8pt
]{principlebox}

% Story callout environment
\newmdenv[
  backgroundcolor=white,
  linecolor=accentgold,
  linewidth=1.5pt,
  topline=false,
  bottomline=false,
  rightline=false,
  innerleftmargin=14pt,
  innerrightmargin=8pt,
  innertopmargin=8pt,
  innerbottommargin=8pt,
  skipabove=10pt,
  skipbelow=6pt
]{storybox}

% TOC
\usepackage{tocloft}
\renewcommand{\cfttoctitlefont}{\Large\bfseries\color{principleblue}}
\setlength{\cftbeforesecskip}{4pt}

% Hyperlinks
\usepackage[
  colorlinks=true,
  linkcolor=principleblue,
  urlcolor=principleblue,
  citecolor=principleblue,
  pdftitle={Enabling the Frontier},
  pdfauthor={Eric Forbes},
  bookmarksnumbered=true
]{hyperref}

% Utilities
\usepackage{enumitem}
\usepackage{graphicx}
\usepackage{booktabs}
\usepackage{xspace}

% Convenient commands
\newcommand{\principle}[1]{%
  \begin{principlebox}
    \textbf{\large #1}
  \end{principlebox}%
}

\newcommand{\sectionintro}[1]{%
  \begin{center}
    \large\itshape\color{midgray} #1
  \end{center}
  \vspace{0.5em}
}

% Editorial placeholder — everything inside a \tbd is scaffolding
% awaiting Eric's input, never book content. Red so it can't be
% mistaken for finished prose when skimming a compiled draft.
\newcommand{\tbd}[1][TBD]{%
  \textcolor{tbdred}{\textit{[#1]}}%
}
